\documentclass[11pt,twoside]{amsart}
\usepackage{amssymb, amsmath, enumerate, palatino, hyperref}
\usepackage[normalem]{ulem}
\usepackage{fullpage}
\usepackage[T1]{fontenc}
\renewcommand{\labelitemi}{\guillemotright}
\usepackage{mathrsfs}

\usepackage{tikz}
\tikzstyle{ball} = [circle,shading=ball, ball color=black,
    minimum size=1mm,inner sep=1.3pt]

\theoremstyle{plain}
\newtheorem{prop}{Proposition}%[section]
\newtheorem{lemma}[prop]{Lemma}
\newtheorem{thm}[prop]{Theorem}
\newtheorem{obs}[prop]{Observation}
\newtheorem{app}[prop]{Application}
\newtheorem*{MainThm}{Main Theorem}
\newtheorem{cor}[prop]{Corollary}
\newtheorem{conj}[prop]{Conjecture}
\theoremstyle{remark}
\newtheorem{rmk}[prop]{Remark}
\newtheorem{prob}{Problem}
\newtheorem{bonus}[prop]{Bonus Problem}
\newtheorem{exc}{Exercise}
\theoremstyle{definition}
\newtheorem{ex}[prop]{Example}
\theoremstyle{definition}
\newtheorem{defn}[prop]{Definition}

\newcommand{\RR}{\mathbb{R}}
\newcommand{\ZZ}{\mathbb{Z}}
\newcommand{\CC}{\mathbb{C}}
\newcommand{\NN}{\mathbb{N}}
\newcommand{\QQ}{\mathbb{Q}}
\newcommand{\PP}{\mathbb{P}}

\newcommand{\cC}{\mathscr{C}}
\newcommand{\Ob}{\operatorname{Ob}}
\newcommand{\Mor}{\operatorname{Mor}}

\newcommand{\Set}{\operatorname{Set}}
\newcommand{\FinSet}{\operatorname{FinSet}}
\newcommand{\Vect}{\operatorname{Vect}}
\newcommand{\FinVect}{\operatorname{FinVect}}
\newcommand{\Mat}{\operatorname{Mat}}
\newcommand{\Gp}{\operatorname{Gp}}
\newcommand{\FinGp}{\operatorname{FinGp}}
\newcommand{\AbGp}{\operatorname{AbGp}}
\newcommand{\Ring}{\operatorname{Ring}}
\newcommand{\CommRing}{\operatorname{CommRing}}
\newcommand{\Field}{\operatorname{Field}}
\newcommand{\Top}{\operatorname{Top}}
\newcommand{\Aut}{\operatorname{Aut}}
\newcommand{\id}{\operatorname{id}}

\newcommand{\defeq}{:=}

\title{Math 113: Discrete Structures\\ Homework 21}
%\author{} %Remove the leading % and put your name here if using this .tex file as a template for your homework.

\begin{document}
\maketitle

\noindent
\textbf{Due:} Friday, March 13 at 10pm.
\bigskip

\begin{prob}\label{Catalan bijections} Consider the full binary tree~$T$:
  \begin{center}
    \begin{tikzpicture}
      \node[ball] at (0,0.2) (0) {};
      \node[ball] at (-1.6,-1) (l) {};
      \node[ball] at (1.6,-1) (r) {};
      \node[ball] at (-2.6,-2) (ll) {};
      \node[ball] at (-0.6,-2) (lr) {};
      \node[ball] at (2.6,-2) (rr) {};
      \node[ball] at (0.6,-2) (rl) {};
      \node[ball] at (-1.6,-3) (lrl) {};
      \node[ball] at (-0.4,-3) (llr) {};
      \node[ball] at (1.6,-3) (rrl) {};
      \node[ball] at (-3.6,-3) (lll) {};
      \draw (0)--(l)--(ll)--(lrl);
      \draw (l)--(lr);
      \draw (rl)--(llr);
      \draw (0)--(r)--(rr);
      \draw (ll)--(lll);
      \draw (r)--(rl)--(rrl);
    \end{tikzpicture}
  \end{center}
  Using the bijections described in our text, find each of the following Catalan
  structures corresponding to~$T$:
  \begin{enumerate}[(a)]
    \item a full parenthesization of the letters~$a,b,c,d,e,f$ (``full'' means
	each multiplication is binary, i.e., involves two factors);
    \item a balanced parenthetical expression with five pairs of~$()$s;
    \item and a Dyck path of length ten.
  \end{enumerate}
  {\bf Note:} It is important to precisely use the conventions used in the
  text.  For example, reordering labels, etc., will not yield the same
  bijection.
\end{prob}

\begin{prob}
  Coin stackings form another set of Catalan structures (i.e., their count is given
  by Catalan numbers).  Here are the~$C_3=5$ coin stackings with a base of three coins:
  \begin{center}
    \begin{tikzpicture}[scale=0.3]
      \draw (0,0) circle (1);
      \draw (2,0) circle (1);
      \draw (4,0) circle (1);
      \begin{scope}[xshift=8cm]
	\draw (0,0) circle (1);
	\draw (2,0) circle (1);
	\draw (4,0) circle (1);
	\draw (1,1.732) circle (1);
      \end{scope}
      \begin{scope}[xshift=16cm]
	\draw (0,0) circle (1);
	\draw (2,0) circle (1);
	\draw (4,0) circle (1);
	\draw (3,1.732) circle (1);
      \end{scope}
      \begin{scope}[xshift=24cm]
	\draw (0,0) circle (1);
	\draw (2,0) circle (1);
	\draw (4,0) circle (1);
	\draw (1,1.732) circle (1);
	\draw (3,1.732) circle (1);
      \end{scope}
      \begin{scope}[xshift=32cm]
	\draw (0,0) circle (1);
	\draw (2,0) circle (1);
	\draw (4,0) circle (1);
	\draw (1,1.732) circle (1);
	\draw (3,1.732) circle (1);
	\draw (2,3.464) circle (1);
      \end{scope}
    \end{tikzpicture}
  \end{center}
  \begin{enumerate}[(a)]
    \item Draw the~$C_4=14$ coin stackings with a base of~$4$ coins.
    \item Prove that the number of coin stackings with a base of~$n$ coins
      is~$C_n$ by describing a bijection between them and Dyck paths of
      length~$2n$. [Hint: consider the region between a Dyck path and the
      diagonal.]
  \end{enumerate}
\end{prob}

\end{document}
